\subsection{半直积}

设 L、N 是群 G 的两个子群，使得 LN = NL；则 LN 是 G 的一个子群，因为

\[
(\mathrm{LN}) (\mathrm{LN}) ^ {- 1} = \mathrm{LNN} ^ {- 1} \mathrm{L} ^ {- 1} = \mathrm{LNL} = \mathrm{LLN} = \mathrm{LN}
\]

此外，映射 \(\varphi: (x, y) \to xy\)（\(N \times L\) 到 \(G\) 中的）是单射，当且仅当 \(N \cap L = \{e\}\)。因为显然若 \(\varphi\) 是单射，则 \(N \cap L = \{e\}\)；反之，关系 \(x' y' = xy\)（其中 \(x, x' \in N\)，\(y, y' \in L\)）蕴含 \(x^{-1}x' = yy'^{-1} \in N \cap L\)；因此若 \(N \cap L = \{e\}\)，则 \(\varphi\) 是单射。于是 \(\varphi\) 是双射，当且仅当 \(LN = G\) 且 \(N \cap L = \{e\}\)。

若 N 是 G 的一个正规子群（或更一般地，在 G 的某个包含 \(N \cup L\) 的子群中是正规的），则条件 LN = NL 自动满足。此外，对每个 \(y \in L\)，映射 \(\sigma_{y}: x \to yxy^{-1}\) 是群 N 的一个自同构，且对 L 的任意两个元素 u、v，我们有

\[
\sigma_ {u v} = \sigma_ {u} \circ \sigma_ {v};\tag{1}
\]

而且对任意 \(x, x'\)（\(\mathbf{N}\) 中的）与任意 \(y, y'\)（\(\mathbf{L}\) 中的），我们还有

\[
(x y) (x ^ {\prime} y ^ {\prime}) = (x \sigma_ {y} (x ^ {\prime})) (y y ^ {\prime}).\tag{2}
\]

反过来：

命题 27 设 N 与 L 是两个群，\(e'\) 与 \(e''\) 是它们各自的单位元。设给定了同态 \(y \to \sigma_y\)（L 到 N 的自同构群 \(\Gamma\) 中的）。则：

\begin{enumerate}
\def\labelenumi{\arabic{enumi})}
\tightlist
\item
  在积集 \(\mathbf{S} = \mathbf{N} \times \mathbf{L}\) 上，内合成法则
\end{enumerate}

\[
(x, y) (x ^ {\prime}, y ^ {\prime}) = (x \sigma_ {y} (x ^ {\prime}), y y ^ {\prime})\tag{3}
\]

定义了一个群结构，对这个结构，\(j_1: x \to (x, e'')\) 是 \(N\) 到 \(S\) 的一个正规子群上的同构，\(j_2: y \to (e', y)\) 是 \(L\) 到 \(S\) 的一个子群上的同构，而 \(\mathrm{pr}_2: S \to L\) 是一个满射同态，其核是 \(j_1(N)\)，并且使得 \(\mathrm{pr}_2 \circ j_2\) 是 \(L\) 的恒等自同构。

\begin{enumerate}
\def\labelenumi{\arabic{enumi})}
\setcounter{enumi}{1}
\tightlist
\item
  设 \(f: \mathbf{N} \to \mathbf{G}\)，\(g: \mathbf{L} \to \mathbf{G}\) 是到一个群 \(\mathbf{G}\) 中的两个同态，使得
\end{enumerate}

\[
f (\sigma_ {y} (x)) = g (y) f (x) g (y ^ {- 1})\tag{4}
\]

对所有 \(x \in \mathbb{N}\) 与所有 \(y \in \mathbb{L}\) 成立。则存在唯一的同态 \(h: \mathbb{S} \to \mathbb{G}\)，使得 \(f = h \circ j_1\) 且 \(g = h \circ j_2\)。

若 \(x, x', x''\) 是 N 的元素，\(y, y', y''\) 是 L 的元素，则

并且

\[
\begin{array}{r l} ((x, y) (x ^ {\prime}, y ^ {\prime})) (x ^ {\prime \prime}, y ^ {\prime \prime}) & = (x \sigma_ {y} (x ^ {\prime}), y y ^ {\prime}) (x ^ {\prime \prime}, y ^ {\prime \prime}) \\ & = (x \sigma_ {y} (x ^ {\prime}) \sigma_ {y y ^ {\prime}} (x ^ {\prime \prime}), y y ^ {\prime} y ^ {\prime \prime}) \\ (x, y) ((x ^ {\prime}, y ^ {\prime}) (x ^ {\prime \prime}, y ^ {\prime \prime})) & = (x, y) (x ^ {\prime} \sigma_ {y ^ {\prime}} (x ^ {\prime \prime}), y ^ {\prime} y ^ {\prime \prime}) \\ & = (x \sigma_ {y} (x ^ {\prime} \sigma_ {y ^ {\prime}} (x ^ {\prime \prime})), y y ^ {\prime} y ^ {\prime \prime}) \end{array}
\]

因而法则 (3) 的结合律由下列事实得出：\(y \to \sigma_{y}\) 是 L 到 \(\Gamma\) 中的一个同态，且 \(\sigma_{y}\) 是 N 的一个自同构。显然 \((e', e'')\) 是 (3) 的单位元，最后

\[
(x, y) \left(\sigma_ {y ^ {- 1}} (x ^ {- 1}), y ^ {- 1}\right) = \left(\sigma_ {y ^ {- 1}} (x ^ {- 1}), y ^ {- 1}\right) (x, y) = \left(e ^ {\prime}, e ^ {\prime \prime}\right)
\]

于是 \((x, y)\) 在 S 中有一个逆元。1) 的其余断言是清楚的。另一方面，由于 \((x, y) = (x, e'') (e', y)\)，满足 2) 的条件的同态 \(h\) 必然满足

\[
h (x, y) = f (x) g (y),
\]

因而若它存在则是唯一的；此外，由 (4) 立即得出

\[
f \left(x \sigma_ {y} (x ^ {\prime})\right) g (y y ^ {\prime}) = f (x) g (y) f (x ^ {\prime}) g (y ^ {- 1}) g (y) g (y ^ {\prime}) = f (x) g (y) f (x ^ {\prime}) g (y ^ {\prime})
\]

这表明 \((x, y) \to f(x)g(y)\) 确是 S 到 G 中的、满足 2) 的条件的一个同态。

推论 命题 27 的 2) 中定义的同态 \(h\) 是单射，当且仅当 \(f\) 与 \(g\) 都是单射且 \(f(\mathbf{N}) \cap g(\mathbf{L}) = \{e\}\)；而 \(h\) 是满射，当且仅当 \(f(\mathbf{N})g(\mathbf{L}) = \mathbf{G}\)。

由于 \(h(x, y) = f(x)g(y)\)，第二个断言是显然的；此外由 (4) 可得 \(f(\mathbf{N})g(\mathbf{L}) = g(\mathbf{L})f(\mathbf{N})\)，而第一个断言由本小节开头的注记得出。

命题 27 中定义的群 S 称为 N 与 L 的外半直积（关于 \(\sigma\) 的）；我们通常把 N（相应地 L）等同于 S 的正规子群 \(j_{1}(N)\){[}相应地子群 \(j_{2}(L)\){]}。若 \(\sigma_{y}\) 是 \(\Gamma\) 的单位元（对所有 \(y \in L\)），我们就重新得到两个群的积的通常概念。

现设 G 是一个群，L 与 N 是 G 的两个子群，使得 LN = NL 且 N 在 NL 中正规，于是对每个 \(y \in L\)，\(\sigma_{y} : x \to yxy^{-1}\) 是 N 的一个自同构，且 \(y \to \sigma_{y}\) 是 L 到 N 的自同构群 \(\Gamma\) 中的一个同态。于是由命题 27 可得：若 S 是 N 与 L 的外半直积（关于 \(\sigma\) 的），则 \(h : (x, y) \to xy\) 是 S 到 G 中的一个同态。h 是双射，当且仅当我们有 \(N \cap L = \{e\}\) 且 NL = G（命题 27 的推论）；这时称 G 为其正规子群 N 与其子群 L 的半直积，且我们常借助于 h 把 G 等同于 S。

命题 28 设 L、N 是两个拓扑群，\(y \to \sigma_y\) 是 L 到 N 的（非拓扑的）群结构的自同构群 \(\Gamma\) 中的一个同态；并设映射 \((x, y) \to \sigma_y(x)\)（\(N \times L\) 到 N 中的）是连续的。则：

\begin{enumerate}
\def\labelenumi{\arabic{enumi})}
\item
  在 N 与 L 的外半直积 S（关于 \(\sigma\) 的）上，N 与 L 的拓扑的积与群结构相容；典范单射 \(j_{1}: \mathrm{N} \to \mathrm{S}\) 与 \(j_{2}: \mathrm{L} \to \mathrm{S}\) 分别是拓扑群 N 与 L 到 S 的子群 \(j_{1}(\mathrm{N})\) 与 \(j_{2}(\mathrm{L})\) 上的同构，而 \(pr_{2}\) 是 S 到 L 上的一个严格态射。
\item
  设 \(f: \mathbf{N} \to \mathbf{G}\) 与 \(g: \mathbf{L} \to \mathbf{G}\) 是到一个拓扑群 \(\mathbf{G}\) 中的两个连续同态，满足 (4)；则同态
\end{enumerate}

\[
(x, y) \rightarrow f (x) g (y)
\]

（S 到 G 中的）是连续的。

这是定义以及积拓扑的性质的一个直接推论。

这样定义的拓扑群 S 称为 N 与 L 的拓扑外半直积（关于 \(\sigma\) 的）；注意，加于 \(\sigma\) 的条件蕴含 L 依照外法则 \((x, y) \to \sigma_{y}(x)\) 在 N 的左边连续作用{[}第 4 目{]}。

现设 G 是一个拓扑群，N 与 L 是 G 的两个子群，使得 G 作为非拓扑群是 N 与 L 的半直积；则显然映射 \((x, y) \to \sigma_{y}(x)\) 在 \(N \times L\) 上连续，且典范双射同态

\[
h: (x, y) \rightarrow x y
\]

（S 到 G 上的）是连续的。但这个同态未必是双连续的；当它双连续时，称 G 为 N 与 L 的拓扑半直积。为此的充要条件是：若 \(p: G \to N\) 与 \(q: G \to L\) 是使 \(z \in G\) 对应于唯一的元素 \(p(z) \in N\) 与 \(q(z) \in L\)（满足 \(z = p(z)q(z)\)）的映射，则映射 p、q 之一是连续的（这时两者都是连续的）。这等同于说，典范映射 \(G \to G/N\) 在 L 上的限制是拓扑群 L 到拓扑群 G/N 上的一个同构。

